% -*- Mode: latex -*- %

\section{Introduction}

We cannot realize the promise of machine learning models---their ability to
ingest huge amounts of data, predict accurately in many domains, their
minimal modeling assumptions---until we can effectively, efficiently, and
responsibly use their predictions in support of scientific inquiry. This paper
takes a step toward that promise by showing how, in problems where
labeled observations or precise measurements of a phenomenon of interest are
scarce, we can leverage predictions from machine-learning algorithms to
supplement smaller real data to validly increase inferential power.

There is no reason to trust black-box machine-learned predictions in place
of actual measurements, calling into question the validity of studies
naively using such predictions in data analyses.
\citeauthor{AngelopoulosBaFaJoZr23}'s
\emph{Prediction Powered Inference} (PPI)~\citep{AngelopoulosBaFaJoZr23}
tackles precisely this issue, providing a framework for valid statistical
inference---computing p-values and confidence intervals---even using
black-box machine-learned models.
%
The idea of PPI is to combine a small amount of labeled data with a larger
collection of predictions to obtain valid and small
confidence intervals.  While the framework applies to a broad range of
estimands and arbitrary machine-learning models, the original
proposal has limitations:
\begin{enumerate}[label=(L\arabic*)]
\item \label{item:intractable}
  for many estimands, such as regression coefficients,
  computing the intervals is intractable, especially in high dimensions;
\item \label{item:adaptivity}
  when the provided predictions are inaccurate, the intervals can
  be worse than the ``classical intervals'' using only the labeled
  data.
\end{enumerate}
\ppipp addresses both limitations.

To set the stage, for observations $(X, Y) \sim \P$ and a loss function
$\loss_\theta(x, y)$, we wish to estimate and perform
inference for the
$d$-dimensional parameter
\begin{equation}
  \label{eq:m-est}
  \theta\opt = \argmin_{\theta \in \R^d}
  \{L(\theta) \defeq \E[\loss_\theta(X, Y)]\}.
\end{equation}
Typical loss functions include the squared error for linear
regression, $\loss_\theta(x, y) = \half(x^\top \theta - y)^2$, or the
negative log-likelihood of a conditional model for $Y$ given $X$,
$\loss_\theta(x, y) = -\log p_\theta(y \mid x)$.  We receive $n$
labeled data points $(X_i, Y_i) \simiid \P$ and $N$ unlabeled data points
$\Xt_i \simiid \Px$, where $X_i$ and $\Xt_i$ have identical distributions.
The key assumption is that we also have access to a black-box model $f$
mapping $X$ to predictions $\what{Y} = f(X)$, which we use to impute
labels for $\Xt_i$. PPI~\cite{AngelopoulosBaFaJoZr23} performs inference on a so-called \emph{rectifier} (an analog of a control variate or debiasing
term) to incorporate these black-box predictions; we reconsider
their methodology to provide two main improvements, adding the
\texttt{\raisebox{.0ex}{++}}.

\paragraph{\texttt{\raisebox{.0ex}{+}}Computational efficiency.}
The original prediction-powered inference constructs a confidence set
for $\theta\opt$ by searching through all possible values of
the estimand $\theta$ and performing a hypothesis test for each.
%
Some estimands, such as means, admit simplifications, but others, such as
logistic regression or other generalized linear model (GLM) coefficients, do
not. When it does not simplify, the original PPI
method~\cite{AngelopoulosBaFaJoZr23} grids the space of $\theta$ and tests
each potential parameter separately; this is both inexact and intractable,
especially in more than two dimensions. 

To address problem~\ref{item:intractable}, we develop computationally fast
convex-optimization-based algorithms for computing prediction-powered point
estimates and intervals for a wide class of estimands, which includes
all GLMs.  The algorithms use the asymptotic normality of these point
estimates and are thus asymptotically exact, computationally lightweight,
and easy to implement.  We show that the resulting intervals are
asymptotically equivalent to an idealized version of the intervals obtained
via the testing strategy (with an infinite grid).  In practice, the
intervals are often tighter, especially in high dimensions.

An important consequence of this theory appears when we wish to estimate
only a single coordinate $\theta\opt_j$ in the model, such as a causal
parameter, while controlling for other covariates and adjusting for nuisance
parameters. Whereas the original PPI~\cite{AngelopoulosBaFaJoZr23} implicitly requires a multiplicity
correction over the entire parameter vector, \ppipp's improved statistical efficiency in the presence
of nuisances allows for inference on single coordinates of regression
coefficients while avoiding multiplicity correction.

\paragraph{\texttt{\raisebox{.0ex}{+}}Power tuning.}
To address problem~\ref{item:adaptivity}, we develop a lightweight
modification of PPI that adapts to the accuracy of the predictive model $f$.
The modification automatically selects a parameter $\lambda$ to interpolate
between classical and prediction-powered inference.  The optimal value of
this $\lambda$ admits a simple plug-in estimate, which we term \emph{power
tuning}, because it optimizes statistical power.  Our experiments and theoretical results indicate that power tuning
is essentially never worse than either classical or prediction-powered
inference, and it can outperform both substantially. Power tuning improves the
performance of the original PPI algorithms when $f$ is not
highly accurate, and, when $f$ is accurate, it can sometimes
guarantee full statistical efficiency.

\section{Preliminaries}

We formalize our problem setting while recapitulating the
prediction-powered-inference (PPI) approach~\cite{AngelopoulosBaFaJoZr23}.
Recall that we have $n$ labeled data points, $(X_i,Y_i) \simiid \P,
i\in[n]$, as well as $N$ unlabeled data points, $\Xt_i \simiid \Px,
i\in[N]$, where the feature distribution is identical across both
datasets, $\P = \Px \times \P_{Y \mid X}$.  In addition to the data, we have
access to a machine-learning model $f$ that predicts outcomes from
features.

Define the population losses
\begin{equation*}
  L(\theta) \defeq \E[\loss_\theta(X, Y)]
  ~~ \mbox{and} ~~
  L^f(\theta) \defeq \E[\loss_\theta(X, f(X))].
\end{equation*}
The starting point of PPI and our approach is the recognition that $\E[\loss_\theta(X, f(X))]
= \E[\loss_\theta(\wt{X}, f(\wt{X}))] = L^f(\theta)$, so that the ``rectified'' loss,
\begin{equation*}
  \LPP(\theta) \defeq L_n(\theta) +
  \widetilde L^f_N(\theta) - L_n^f(\theta),
\end{equation*}
where
\begin{equation*}
 L_n(\theta) \defeq \frac{1}{n} \sum_{i = 1}^n \loss_\theta(X_i, Y_i), ~~ L_n^f(\theta) \defeq \frac{1}{n} \sum_{i = 1}^n \loss_\theta(X_i, f(X_i))
  ~~ \mbox{and} ~~
  \wt{L}_N^f(\theta) \defeq \frac{1}{N} \sum_{i = 1}^N \loss_\theta(\wt{X}_i,
  f(\wt{X}_i)),
\end{equation*}
is unbiased for the true objective: $\E[\LPP(\theta)] = L(\theta)$.

%Prediction-powered inference (PPI)~\cite{AngelopoulosBaFaJoZr23} combines
%the labeled data and the unlabeled data to perform more powerful inferences
%than classical procedures using only the labeled data in an assumption-light
%way, requiring only that $X$ and $\Xt$ have identical marginal
%distributions, treating $f$ as a black-box. 
%%This avoids the assumption- and
%%computation-heavy approaches that semiparametric inference typically
%%requires~\cite{RobinsRoZh94, BickelKlRiWe98, Imbens04, ImbensRu15,
%%  ChernozhukovChDeDuHaNeRo16}.
In the original, PPI first forms
$(1-\alpha)$-confidence sets for the gradient of the loss $\nabla L(\theta)$, denoted $\C^{\nabla}_\alpha(\theta)$, 
and then constructs the confidence set $\CPP_\alpha$ for $\theta\opt$ as
\begin{equation}
  \label{eq:CPP-original}
  \CPP_\alpha = \{\theta \in \R^d ~ \mbox{s.t.} ~ \zeros
  \in \C^\nabla_{\alpha}(\theta)\}.
\end{equation}
Constructing $C^\nabla_\alpha(\theta)$ for a \emph{fixed} $\theta$ is a trivial application of the central limit theorem to $\nabla \LPP (\theta)$. 
The issue is that forming the confidence set for \emph{every} $\theta \in \R^d$ is computationally infeasible, because one cannot simply iterate through every possible value.

%Another problem with the construction \eqref{eq:CPP-original} is that $\CPP_\alpha$ is a confidence region for \emph{all} $d$ coordinates of $\theta$; this multiplicity is intrinsic and cannot be avoided straightforwardly.
% If we are only interested in a single coordinate and the other coordinates are nuisance parameters---such as in the case of inference on a causal effect in the presence of control variables---this type of simultaneous coverage is overly conservative.
%\citet{AngelopoulosBaFaJoZr23} use the
%gradient of the
%``rectified'' loss
%\begin{equation*}
%  \LPP(\theta) \defeq L_n(\theta) +
%  \widetilde L^f_N(\theta) - L_n^f(\theta),
%\end{equation*}
%which is evidently unbiased for the true objective $L(\theta) =
%\E[\LPP(\theta)]$, and the central limit theorem (or, to guarantee finite
%sample validity, concentration inequalities). Using that
%$G_n(\theta) \defeq \nabla \LPP(\theta)$ satisfies $\E[G_n(\theta)] = \nabla
%L(\theta)$, whenever $n/N \to r$ then
%\begin{equation*}
%  \sqrt{n} (G_n(\theta) - \nabla L(\theta))
%  \cd \normal(0, \Sigma)
%  ~~ \mbox{for} ~~
%  \Sigma = r \cov(\nabla \loss_\theta(X, f(X)))
%  + \cov(\nabla \loss_\theta(X, Y) - \nabla \loss_\theta(X, f(X))).
%\end{equation*}
%\citeauthor{AngelopoulosBaFaJoZr23} take
%take consistent estimates $\what{\sigma}_j^2$ of the diagonal $\Sigma_{jj}$
%and set
%\begin{equation}
%  \label{eqn:bonferdingdong}
%  \C^\nabla_\alpha(\theta) \defeq
%  \left\{G_n(\theta) + v ~ \mbox{s.t.} ~
%  v \in \R^d ~ \mbox{satisfies} ~ |v_j| \le
%  \frac{z_{1 - \alpha / (2d)}}{\sqrt{n}}
%  \cdot \what{\sigma}_j ~ \mbox{for~} j \in [d]\right\},
%\end{equation}
%where $z_{1 - \alpha}$ denotes the $(1 - \alpha)$-quantile of a standard
%normal.  That this is asymptotically valid is
%immediate~\cite[Thm.~B.1]{AngelopoulosBaFaJoZr23} as $\nabla L(\theta\opt) =
%\zeros$, so $\liminf_n \P(\theta\opt \in \CPP_\alpha) = \liminf_n
%\P(\zeros \in \C^\nabla_\alpha(\theta\opt)) \ge 1 - \alpha$.
%% In Section~\ref{sec:equivalence}, we revisit the
%% construction~\eqref{eq:CPP-original}, providing a tighter version
%% of it that avoids the multiplicity penalties the (naive)
%% confidence sets~\eqref{eqn:bonferdingdong} imply.


\subsection{Computational efficiency: overview}

We take PPI as our point of departure, but adopt a different perspective. In particular, the
rectified loss leads to the \emph{prediction-powered point estimate},
\begin{align}
  \label{eq:ppi}
  \tag{$P_1$}
  \thetaPP &=\argmin_\theta \LPP(\theta), \text{ where } \LPP(\theta) \defeq L_n(\theta) + \widetilde L^f_N(\theta)  - L_n^f(\theta).
\end{align}
This estimate (along with its forthcoming variant~\eqref{eq:ppi-lambda}) is the
main object of study in this paper, and we derive its asymptotic normality
about $\theta\opt$.  This result is the keystone to our efficient
\ppipp algorithms, allowing us to construct confidence intervals of the form
\begin{equation}
  \label{eq:CPP-new}
  \CPP_{\alpha} = \left\{\thetaPP_j \pm z_{1-\alpha/2} \cdot \hat\sigma_j / \sqrt{n} \right\},
\end{equation}
where $\hat\sigma_j^2$ is a consistent estimate of the
asymptotic variance of $\thetaPP_j$ and $z_{q}$ denotes the
$q$-quantile of the standard normal distribution. Similar confidence intervals centered around $\thetaPP$ apply when inferring the whole vector $\theta\opt$ rather than a single coordinate $\theta\opt_j$.

It should be clear that the set~\eqref{eq:CPP-new} is easier to compute than the set~\eqref{eq:CPP-original}.
The latter requires the analyst to ``iterate through all $\theta \in \R^d$'' and test
optimality separately; the former has no such requirement, as \ppipp
simply requires placing error bars around a point estimate.  
This perspective, though simple, results in far more efficient algorithms.
In Section~\ref{sec:ppi++_glms}, we provide algorithms for generalized linear models, and in Section~\ref{sec:ppipp-m-estimation}, we show the case of general M-estimators. 
In Section~\ref{sec:equivalence}, we show that the \ppipp intervals are asymptotically equivalent to the naive strategy~\eqref{eq:CPP-original}, so the increased computational efficiency comes with no cost in terms of statistical power.

\subsection{Power tuning: overview}

The prediction-powered estimator~\eqref{eq:ppi} need not outperform classical inference when $f$ is inaccuate:
limitation~\ref{item:adaptivity} of PPI. 
We thus generalize PPI to allow it to adapt to the ``usefulness'' of the supplied predictions, yielding an estimator never worse than classical inference. 
To do so, we introduce a tuning parameter $\lambda \in \R$, defining
\begin{equation}
\label{eq:ppi-lambda}
\tag{$P_\lambda$}
\thetaPP_\lambda = \argmin_\theta \LPP_\lambda(\theta), \text{ where } \LPP_\lambda(\theta)\defeq L_n(\theta) + \lambda \cdot (\widetilde L^f_N(\theta) - L_n^f(\theta)).
\end{equation}
Certainly the objective remains unbiased for any $\lambda$, as
$\E[\LPP_\lambda(\theta)] = L(\theta)$. The problem \eqref{eq:ppi-lambda}
recovers \eqref{eq:ppi} when $\lambda = 1$ and reduces to performing
classical inference, ignoring the predictions, when
$\lambda=0$. 
In Section~\ref{sec:setting-lambda}, we show how to optimally tune $\lambda$ to maximize estimation and inferential efficiency, yielding a data-dependent tuning parameter $\hat{\lambda}$.
This adaptively chosen $\hat{\lambda}$ leads to better estimation and inference than both the classical and PPI strategies.
(In fact, sometimes it is fully statistically efficient, though that is
not our focus here.)
Section~\ref{sec:experiments} puts everything together through a series of experiments highlighting the benefits of \ppipp.


\subsection{Related work}

\citet{AngelopoulosBaFaJoZr23} develop prediction-powered inference (PPI)
to allow estimation and inference using black-box machine learning models;
our work's debt to the original proposal and motivation is clear.
%% %
%% The idea is to use the predictions in combination with a (usually small)
%% labled dataset to achieve more powerful inferences than possible with the
%% labeled dataset alone.
%% %
%% The main difference between the original PPI proposal and \ppipp is that our
%% paper explicitly derives the asymptotic distribution of the
%% prediction-powered point estimator, allowing for efficient confidence
%% interval constructions.
%% %
%% Our approach also introduces power tuning, which significantly tightens the
%% intervals while ensuring that PPI always outperforms the classical approach.
%
The PPI observation model---where we receive $n$ labeled observations and
$N$ unlabeled observations---also centrally motivates semi-supervised
learning and inference~\cite{ZhangBrCa19, ZhangBr22, SongLiZh23}.
%
This line of work also considers the problem of inference with a labeled
dataset and an unlabeled dataset; the key point of departure is PPI's access
to a pretrained machine learning model $f$.
%
Fitting empirically strong predictive models of a target $Y$ given $X$ has
been and remains a major focus of the machine learning
literature~\cite{LeCunBeHi15}, so that leveraging these
models as black boxes, in spite of the lack of essentially \emph{any}
statistical control over them (or even knowledge of how they were fit),
provides a major opportunity to leverage modern prediction techniques.
%
This enables PPI to apply to estimation problems where, for example,
high-performance deep learning models exist to impute labels (such as
AlphaFold~\cite{JumperEvEtAl21} for biological structures or
CLIP~\cite{RadfordKiHaRaGoAgSaAsMiClKrSu21} for image labeling), and, though we may have no
hope of understanding precisely what the models are doing, to obtain more
powerful inferences in those settings.

\newcommand{\thetaAIPW}{\hat{\theta}^{\textup{AIPW}}}

One may also take a perspective on the \ppipp approach coming from the
literatures on semiparametric inference, causality, and missing
data~\cite{RobinsRo95, BickelKlRiWe98, RobinsRoZh94, Tsiatis06,
  ChernozhukovChDeDuHaNeRo16}, as well as that on survey
sampling~\cite{SarndalSwWr03}.  We only briefly develop the connections, but
it is instructive to formulate our problem as one with missing data, where
the dataset $\{\Xt_i\}_{i=1}^N$ has no labels. Consider now augmented
inverse propensity weighting (AIPW) estimators~\cite{RobinsRo95, Tsiatis06},
focusing in particular on estimating the mean $\E[Y]$. Then because we know
the ratio $r = n/N$ (the propensity score, or probability of missingness), AIPW~\cite{RobinsRo95, Tsiatis06}
estimates $\E[Y]$ by
\begin{align}
  \thetaAIPW & \defeq \frac{1}{n} \sum_{i = 1}^n
  (Y_i - \hat{\E}[Y \mid X = X_i])
  + \frac{1}{N + n} \bigg(\sum_{i = 1}^n \hat{\E}[Y \mid X = X_i]
  + \sum_{i = 1}^N \hat{\E}[Y \mid X = \wt{X}_i]\bigg) \nonumber \\
  & \, = \frac{1}{n} \sum_{i = 1}^n
  (Y_i - f(X_i))
  + \frac{1}{N + n} \bigg(\sum_{i = 1}^n f(X_i)
  + \sum_{i = 1}^N f(\wt{X}_i)\bigg)   \label{eqn:aipw}
  \\
  & \, = \frac{1}{n} \sum_{i = 1}^n
  Y_i 
  + \frac{1}{1 + \frac n N} \bigg(-\frac 1 n \sum_{i = 1}^n f(X_i)
  + \frac 1 N \sum_{i = 1}^N f(\wt{X}_i)\bigg),
  \nonumber
\end{align}
where in the first line $\hat{\E}$ denotes the estimated conditional
expectation of $Y$ given $X$, and in the second we replace it with the
black-box predictor $f$.
%
This corresponds to the particular choice $\lambda = \frac{1}{1 + r}$ in the
\ppipp estimator~\eqref{eq:ppi-lambda}.
%
The point of emphasis in PPI is distinct, however; while semiparametric
inference in principle allows machine learning to estimate nuisance
parameters ($\E[Y \mid X]$ in the AIPW estimator~\eqref{eqn:aipw}), we
develop easy-to-compute, practical, and assumption-light estimators, where
we truly treat $f$ as an unknown black box, and \ppipp leverages the
predictions as much as is possible.
%
The practicality is essential (as our experiments highlight), and it
reflects the growing trend in science to use the strength of modern
machine-learned predictors $f$ to impute labels $Y$.
%
The debiasing~\eqref{eq:ppi-lambda} thus directly allows arbitrary
imputation of missing labels, avoiding challenges in more common imputation
approaches~\cite{Rubin18}.
%
We remark in passing that survey sampling approaches frequently use
auxiliary data to improve estimator efficiency~\cite{SarndalSwWr03}, but
these techniques are limited to simple estimands and do not use machine
learning.
%

Our work can be seen as an improvement of PPI along the lines of Veridical Data Science~\cite{Yu18, YuKumbier20, YuBa24}, the main tenets of which are predictability, computability, and stability (PCS).
%
On the note of predictability, PPI can diminish statistical power if the model $f$ cannot adequately predict $Y$; our procedure will automatically check if this is the case and reduces to classical inference if the model is bad. Therefore, if classical inference is believed to yield statistically valid results, \ppipp will be valid as well.
%
On the note of computability, for algorithms like logistic regression, PPI intervals were not efficiently computable, while \ppipp intervals are.


See also \citet{AngelopoulosBaFaJoZr23} for more discussion of related work.

%% More distantly, PPI is also related to the literatures on semiparametric inference~\cite{RobinsRo95, BickelKlRiWe98}, survey sampling~\cite{SarndalSwWr03}, and inference with missing data~\cite{RobinsRoZh94}.
%% Semiparametric inference also considers the use of black-box machine learning models to estimate a nuisance parameter, but does not usually consider the semi-supervised setting.
%% The survey sampling literature uses of auxiliary data to improve the efficiency of estimators, but these techniques are limited to a small set of simple estimands and do not use machine learning.
%% Our work can also be seen as a missing data problem where the labels are missing on the unlabeled dataset.
%% A common solution to the missing data problem is to impute the missing labels via prediction algorithms~\cite{Rubin18}. 
%% However, much of the literature on imputation does not have a principled way of debiasing the prediction errors for the purpose of computing downstream estimands. 
%% For a more detailed related owrk section, see~\cite{AngelopoulosBaFaJoZr23}.
